\section{Discussion}

Our experiments validate that task-conditioned conversion preserves adaptation capability in sub-quadratic architectures.

\subsection{Key Findings}

\textbf{The efficiency-adaptation tradeoff is not fundamental.} TC-SSM achieves 0.25\% accuracy gap with 0.85$\times$ overhead---both better than thresholds set for acceptable tradeoffs. This suggests that the apparent conflict between sub-quadratic efficiency and adaptation capability stems from how conversion is performed, not from architectural constraints.

\textbf{Self-supervised task discovery works.} Cluster purity of 0.74 confirms that meaningful task structure exists in pretrained hidden states. This enables TC-SSM to condition on task identity without requiring task labels during conversion.

\textbf{Task conditioning can be computationally free.} The sub-unity overhead (0.85$\times$) was unexpected. We hypothesize this results from improved memory bandwidth utilization: task embeddings are small (32-dimensional), cached, and reused across all sequence positions.

\subsection{Limitations}

\textbf{Single model scale.} All experiments use BERT-base (110M) to Mamba-130M conversion. We cannot claim results generalize to larger scales (1B+), where different phenomena may emerge.

\textbf{SuperGLUE classification only.} Our evaluation covers NLU classification tasks. Generation, reasoning, and multimodal tasks remain untested.

\textbf{Sample imbalance effects.} Embedding quality correlates with training data size. Small tasks (CB: 56 samples, COPA: 100) achieve 0\% individual linear probe accuracy despite contributing to overall discriminability.

\subsection{Broader Impact}

TC-SSM enables efficient deployment of adaptable models, potentially democratizing access to capable NLP systems on resource-constrained hardware. However, more efficient models may accelerate deployment in contexts where careful oversight is warranted. We recommend evaluation protocols that assess behavior across diverse task types before deployment.
